\section{Remaining gates, false formulations, and the exact delta}
\label{sec:r11-remaining-gates}

The status of a universal assertion is not changed by
proving a qualified theorem with a specific boundary
class.  Conversely, a counterexample to an unqualified
assertion should not be recorded merely as an open
problem.  The following ledger distinguishes the two.

\subsection{The capacity and unbounded-calculus boundary}

\paragraph{Proved unbounded sector.}
The original arbitrary-background construction now includes
coefficient-finite non-polar closedness on ordinary germs;
compact-sublevel outer quasicontinuity; the exact no-loss
capacity-tail cone; its comparison, domination and
relative uniqueness; all positive power-energy classes;
all-degree curvature recursion and graded first-order
and energy currents; and their precise capacity and
strong-energy continuity statements.
On one original affinoid neighborhood, the new analytic
tester extension compares the full local capacities
on compact cores.
On strict analytic boxes with an extending background,
the finite-mass boundary tester core computes raw
capacity, and the full global $E^1$ Green and energy
geometry hold.
These are \texttt{NEW\_PROVED} results with the
individual assumptions printed in their modules.

\paragraph{General raw and outer capacities: \texttt{OPEN}.}
An unrestricted proof of
$C_S(K;V)=C_S^*(K;V)$ for every compact $K$,
or of raw/outer equality on every Borel set, is not
provided.  The elementary continuity of the outer
regularization on decreasing compact sets was already
frozen and is not the missing assertion.
The remaining weak-limit gate is concrete: for a
locally weakly convergent family of full normalized
original tester measures, one needs control on a
specified compact or Borel set by the raw tester
capacity, rather than merely the outer-capacity
bound inherited on opens.
No infimum over neighborhoods has been interchanged
with a supremum over testers.

The exact raw equality for the finite-mass core
(TC6) concerns two tester families on the same
domain.  The domain comparisons (AL8)--(AL9)
compare two capacities of the same kind.
The strict-gap equality (BC17) and the scalar
coarea estimates concern specified sets.
None of these statements proves the unrestricted
raw/outer assertion.

\paragraph{A single auxiliary energy tail: \texttt{FALSE}.}
The new fixed-tail example (FT2) compares the full
nonlinear capacity of one background with the
linear energy capacity of a particular auxiliary
tail: its point has positive nonlinear capacity
and zero auxiliary energy capacity. The frozen R6
example instead compares the BT and energy capacities
of the same degree-one background and gives a BT-null
compact set of positive energy capacity.
Together they explain why neither identifying
the two polar ideals nor controlling an entire
nonlinear tester family by one fixed linear
form resolves the raw-capacity gate.
They are not counterexamples to raw/outer
regularity itself.

\paragraph{Unrestricted recursion: \texttt{FALSE}.}
The logarithmic example (CR2) has a locally
coefficient-integrable multiplier
$uN_{q-1}$, yet
$D(uN_{q-1})=\delta_o$ while the non-polar
top product is zero.
No-loss and the printed integrability
condition repair that passage; every positive
power class supplies sufficient integrability.
The discarded polar residue cannot be removed
by redefining the measure on its pole after
differentiation.

\paragraph{Unrestricted endpoint continuity: \texttt{FALSE}.}
Uniform $E^\sigma$ energy and capacity convergence
alone do not give critical weighted continuity.
The shell family in (PS7) has a nonzero limiting
weighted atom.
The proved repairs are fixed-tuple tails,
uniform capacity-tail integrability, actual
decreasing lower barriers, or the strong
$E^1$ discrepancy topology with its derived
uniform tails.

\subsection{Poles, envelopes, and admissible representatives}

\paragraph{Complete compact-sublevel pole sets.}
The logarithmic construction gives one actual
representative of each such pole set in every
positive power class, simultaneously for all
allowed extending backgrounds.
This is \texttt{NEW\_PROVED}.
The ambient R2 pole/negligible statements on
their specified exhaustion domains remain
\texttt{INTERNAL\_INPUT}.
The unrestricted claim that a relative
capacity-null set must be an ambient or
compact-sublevel complete pole set is
\texttt{FALSE} without nondegeneracy and
domain qualifications.
An appropriate converse with explicit
nondegeneracy, representation and domain
hypotheses remains \texttt{OPEN}.

\paragraph{Ordinary upper regularization: \texttt{FALSE}.}
The all-power boundary-core example (PO1)--(PO7)
retains a common actual lower barrier, finite
mass, unit energies and even an actual
classical solution of the prescribed source.
Its ordinary upper-regularized Perron
supremum nevertheless loses the source.
This is a strengthening of the frozen
proper-support example, not a new proof
of the original bounded failure.

\paragraph{Relative completed envelopes: \texttt{NEW\_PROVED}.}
The barrier completion constructs original
closed currents, derivatives and allowed
energy currents before asserting their
unbounded recursive identities.
Its unbounded relative Perron envelope
preserves Borel constraints and a common
source lower bound.
The completed barrier-constrained
variational problem has a unique relative
maximizer, a precise Euler inequality,
optimal weighted-data stability, and a
retraction on its finite-mass operator image.
These are actual existence statements in
an explicitly enlarged Borel potential
class.  They do not claim that every
completed potential has an ambient
classical representative.

\paragraph{Full relative strong completion: \texttt{NEW\_PROVED}.}
The first-energy strong space removes the prescribed actual
lower barrier. Uniform deep tails and global weighted tightness
give its complete actual-Borel/current realization. Its full
comparison, domination, strict locality, exact truncation
carriers and finite-mass approximation retain the encoded
strong boundary class. The intrinsic capacity-layer formula
characterizes each power moment of an already constructed
completed member; its strict-gap raw/outer equality still
does not give unrestricted capacity regularity.

The exact intrinsic characterization now requires same-domain
bounded capacitary admissibility of every truncation and a
finite outer first-energy capacity integral. This removes
the need to prescribe a strongly Cauchy defining row, while
retaining a genuine analytic admissibility gate. It also
gives the precise monotone closure theorem. Every nonempty
strong completed family has a relative Perron envelope by
energy minimization over finite maxima, preserving source
lower bounds and Borel upper constraints.

For positive data dominated by a constant times one actual
first-energy top measure, the full strong cone has a unique
variational maximizer, sharp energy bounds and quantitative
data stability. Total source mass at the boundary may be
infinite. Literal datum integrals are proved continuous from
the completed tester estimate before maximization is used.
The source equation remains distinct even in this explicit
dominated class.
Within one fixed dominated interval, the actual completed
operator image is closed for the proved absolute
energy-dual distance. This is an approximation criterion:
it requires approximants already in that operator image
and does not assert that the image fills the interval.

\paragraph{Maximality implies the source equation: \texttt{FALSE}.}
The completed maximizer in (BV13), and even the
unconstrained strong-cone maximizer in (VF13), can
fail the prescribed current equation.
The R10 compact-source rigidity extends through
the new degree-zero current convergence to the
entire strong completion: a uniform first-energy
bound controls local convex Gauss traces by (SO3).
Thus a classical representative is not
the only missing issue: the constrained
admissible variation cone is itself a
substantive equation gate.
The exact proved condition is (BV8).
The current equation follows only under
additional valid variations, as explicitly
stated in the actual-class conditional
corollary.

\subsection{Dirichlet and obstacle problems}

\paragraph{The complete relative linear measure-data problem.}
For an arbitrary original degree-$(n-1,n-1)$
background on the printed affinoid part domain,
every finite signed energy-diffuse Radon datum
has a unique relative renormalized zero-boundary
solution.  Its original current equation,
exact level carriers, bounded-value coarea,
hard-curvature global total-variation convergence,
and reference independence are
\texttt{NEW\_PROVED}.
The signed finite-diffuse obstacle problem
likewise has a unique least majorant and
an exact positive contact reaction, under
the stated finite-curvature obstacle
assumption.
The reaction bound and source stability
are proved before and during the limiting
construction; contact is not inferred
from an uncontrolled weak restriction.

The word zero boundary here denotes the
specified part-space clipping and radiation
conditions.  It does not denote arbitrary
ordinary pointwise boundary values.
Unweighted energy of a measure-data
solution is not always Radon locally;
the explicit positive-data example
proves this assertion \texttt{FALSE}.
The chain and coarea theorems use only
the finite bounded-value energy measures
they actually require.

\paragraph{The nonlinear equation remains appropriately qualified.}
For $q>1$, the new calculus identifies the
derivative of every actual admissible
power-class segment and solves its
relative linearized signed diffuse
problem on a small part domain.
This is a \texttt{DIRECT\_COROLLARY};
it does not realize an arbitrary energy
tangent as a classical variation.

An unrestricted assertion that every
capacity-controlled or finite-dual-energy
datum has an ambient classical solution
for every original $S$ is already
\texttt{FALSE} in the precise remaining-degree-one
moment example, now also in its unbounded
power classes and full strong completion.
The genuinely open nonlinear program is
therefore to specify a valid admissible
data class, potential class and boundary
condition, and prove existence, uniqueness,
comparison and suitable regularity there.
No theorem here removes those three
specifications.

\paragraph{A precise degree-one image criterion.}
For $q=1$ the strong first-energy completion embeds
isometrically as the closed cone $\mathcal C$ in the
original part Hilbert space. The full finite-$H$-dual
source equation is solvable in that cone exactly when
its Riesz vector lies in $\mathcal C$, equivalently
when the stated polar-cone inequalities hold.
The source-functional projection, its sharp stability,
and its exact polar defect are direct Hilbert consequences.
The original smooth core is a difference of actual
finite-mass common-collar potentials, so the cone has
dense linear span. Its properness shows why dense span
does not provide feasible two-sided directions at a
specified maximizer.
This is an exact criterion in the enlarged relative
class, not a constructive classification of all ambient
classical data or an identification of the BT and energy
polar ideals. Pairings with general signed $H$-dual data
remain continuous energy functionals unless their literal
measure interpretation is separately justified.

\subsection{The analytic relation quotient is still separate}

\paragraph{Full GQ: \texttt{OPEN}.}
For every compatible finite actual PS/PL
presentation $e$, the target is still
\[
 e=0\text{ as an actual analytic delta-form}
           \quad\Longrightarrow\quad A_S(e)=0
 \tag{OGQ}
\]
for every fixed original arbitrary $S$,
on every ordinary germ, after all actual
moment refinements, and against all original
complementary tests.
The positive-zero kernel, fixed-moment
top-real-degree signed relations, and the
R8--R10 matroidal, fan and actual analytic
sectors remain frozen inputs.
The new quotient of actual power-class
potentials modulo relative capacity
does not imply (OGQ).
No equality on the full GK delta-test
category has been silently added.

\paragraph{Low-real-degree CTO: \texttt{OPEN}.}
Signed constant-tensor order transfer in
$1<b<r$, including the critical
$b=2<r=3$ range, is not completed by
the new ordinary energy and capacity
identities.  They do not provide an
arbitrary positive lift of an actual
signed analytic relation.
The precise R10 conormal and order
sectors remain as frozen, with their
counterexamples to stronger assertions.

\paragraph{Other frozen open gates.}
$W_{\rm loc}$ density and the Gauss-residue
auxiliary-field quotient remain
\texttt{OPEN}.  They are not inputs to
any R11 new proof.
The old internal R/B/A2 results remain
authorized \texttt{INTERNAL\_INPUT},
and completed A2 geometric adelic
nefness remains completed.

\subsection{Verification, provenance and continuation}

Every module is retained as exact source,
with independent review records binding
the reviewed SHA256.
The standalone manuscript is generated
from those modules with a recorded source
map; its layout substitutions do not
change mathematics.
The exact midpoint script is an auxiliary
rational normalization check only.
The infinite-degree proof is the finite
binomial argument in (SM3)--(SM7), not
the finite range checked by the script.

Two provenance corrections are explicit.
The R6 polar-ideal separation was
rediscovered independently and excluded
from new-proof counts; its complete
draft remains in the failed-attempt record.
The bounded top weak maximum inequality
was already frozen in R3; R11 adds its
exact all-degree positive level-current
decomposition and the unbounded extension.
The R3/R4 bounded midpoint and constrained
retraction mechanisms are also credited;
R11 supplies the global unbounded domain
and the new sharp constant.

The minimal next tasks are mathematical:
raw weak-limit domination for the full
normalized tester family; a valid source
or variation criterion beyond the now
explicit false unrestricted one; completed
nonlinear contact/equation mechanisms with
their true boundary class; and the
independent full GQ/low-degree CTO gates.
The completed barrier calculus and
finite signed diffuse contact theorem
can be used as new inputs, with all their
domain and representative assumptions
carried forward.

\paragraph{External source used for the new measure-data extension.}
T.~Klimsiak and A.~Rozkosz,
\emph{On the structure of bounded smooth measures
associated with a quasi-regular Dirichlet form},
arXiv:1410.4927v2, Lemma 2.1:
\url{https://arxiv.org/abs/1410.4927}.
The exact finite-zero-energy nest input and
its hypotheses are stated in the main
Dirichlet proof.  General form-theoretic
facts already used in the frozen faithful
energy theory retain that earlier dependency
record; this round does not add an
unidentified PDE existence theorem.
